\chapter{Discussão integrada}
\label{chap:discussao}

\section{Três operações com consequências distintas}

A inferência da massa do gráviton depende tanto da resposta dispersiva quanto da distribuição dos estimadores usados para medi-la. Os experimentos distinguem três operações: aproximar a distribuição de formas quadráticas por uma normal, comprimir a informação em frequência e substituir a resposta dispersiva por uma resposta de referência. A separação dessas operações é o principal resultado metodológico da pesquisa.

A0 conserva os coeficientes complexos do experimento periódico. A usa estimadores quadráticos e uma verossimilhança normal com seus momentos. B comprime A por uma transformação linear fixa; C analisa esse mesmo dado comprimido com uma resposta aproximada. O controle gaussiano G permite comparar A e B com distribuições corretamente especificadas, isolando a perda de informação associada à compressão.

\section{Quando os momentos não determinam a verossimilhança}

As falhas de calibração das aproximações normais são um resultado científico central. Mesmo com média e covariância corretas, uma forma quadrática de variáveis gaussianas pode apresentar assimetria e caudas incompatíveis com a normal. As dez rejeições da campanha tensorial que se mantêm ao variar os resultados dentro de sua incerteza numérica demonstram essa insuficiência nos cenários examinados. Os estudos escalar e de população ampliada fornecem exemplos adicionais, com desenhos experimentais próprios.

A cobertura central nominal de 90\% para $\log_{10}\mathrm{EFAC}$ foi de 313/500 (62,6\%) em A\_CN e 412/500 (82,4\%) em B\_CN. Em A\_CN, mesmo o intervalo de contagens permitido pela incerteza numérica, de 299 a 325 inclusões, permanece distante das 450 inclusões esperadas nominalmente.

Essa conclusão muda a interpretação dos contrastes entre respostas. Em B\_CN e C\_CN, ambos os ajustes usam uma distribuição aproximada para os estimadores físicos. Seu contraste mede a mudança de resposta dentro dessa mesma aproximação; um pequeno deslocamento entre posteriores pode coexistir com cobertura inadequada em ambas. A calibração de B é, portanto, parte necessária da avaliação de C. Já o contraste A\_G/B\_G quantifica a compressão em um experimento gaussiano corretamente especificado.

Densidades de formas quadráticas ou aproximações que incorporem cumulantes superiores são caminhos para melhorar o modelo probabilístico. Uma expansão de Edgeworth, por exemplo, também precisaria ser verificada quanto à normalização, positividade e precisão nas caudas relevantes. Os resultados estabelecem a necessidade de avaliar a adequação da densidade nos regimes estudados; a escolha de uma correção depende desse exame \autocite{franciolini2025}.

\section{Limite superior, informação e dependência da priori}

Uma posterior calibrada pode conter pouca informação sobre a massa. O exemplo mais direto é uma verossimilhança constante com priori uniforme em $u\in[0,1]$: a posterior coincide com a priori, $D_{\rm KL}=0$ e $Q_{0,95}(u)=0,95$. Como $u=f_gT$, obtém-se
\begin{equation}
 Q_{0,95}(m_gc^2)=0,95\,\frac{h_{\mathrm P}}{T}.
 \label{eq:discussao-priori}
\end{equation}
O valor numérico é determinado pelo suporte escolhido. A proximidade de um limite com essa escala motiva comparar explicitamente priori e posterior; isoladamente, ela tampouco diagnostica a origem de um limite publicado. Para buscas de dispersão em PTAs, a apresentação conjunta de limites, ganho de informação e sensibilidade à priori esclarece quanto da restrição vem dos dados.

No controle gaussiano da população ampliada, a diferença média entre as divergências posterior--priori de A\_G e B\_G foi $0,06211\pm0,00477$ nat em P12K4 e $0,10631\pm0,00695$ nat em P16K8. Os erros são erros Monte Carlo da média sobre 500 realizações da priori. Essas diferenças estimam a perda de informação esperada da compressão; sua ordenação é uma propriedade da média populacional, e pode variar entre realizações.

A precisão dos contrastes B/C depende também da dimensão da inferência. Com marginalização dos parâmetros auxiliares, os contrastes primários permaneceram inconclusivos na escala de 0,01 da largura da priori. No experimento complementar com ruído fixado e 50 realizações, os intervalos numéricos das quatro diferenças médias de $Q_{95}(u)$ ficaram dentro de $[-0{,}01;0{,}01]$ (Seção~\ref{sec:contrastes-condicionais}). Esse resultado estabelece pequenos deslocamentos médios na amostra condicional e delimita a pergunta ainda aberta: como a marginalização modifica essa comparação.

\section{Polarizações vinculadas e degenerescência no limiar}

A reconstrução do TG conecta a curvatura local à propagação massiva, enquanto Fierz--Pauli define o vínculo escalar usado na extensão estatística. As deformações transversal e longitudinal desse modo pertencem à mesma amplitude. A relação $p_l=-\chi p_b$, com $\chi=(f_g/f)^2$, e os termos cruzados da resposta são essenciais para preservar essa estrutura física.

No limiar exato, $\Gamma_0=\Gamma_T/2$. Em um canal e sob a mesma geometria, a correlação depende apenas de $S_T+S_0/2$: a forma angular, sozinha, é incapaz de separar as duas potências. Esse resultado interessa às buscas de polarizações não tensoriais, pois diferentes conteúdos de polarização podem produzir a mesma correlação angular nesse limite. Perto do limiar, a proporcionalidade é aproximada e sua separação depende da sensibilidade; em uma banda extensa, a evolução em frequência pode quebrar a degenerescência. A conclusão se refere ao modo vinculado de Fierz--Pauli e à normalização adotada, e não a toda população escalar possível.

A análise de Fisher caracteriza essa degenerescência local. Sua projeção sobre parâmetros auxiliares descreve direções de pouca informação na vizinhança de um ponto, complementando a recuperação posterior em experimentos com ruído conhecido.

\section{Relação com a literatura}

Han e Zhao investigam previsões de massa com covariância completa de estimadores PTA--astrometria \autocite{han2026}; Franciolini e colaboradores examinam a adequação de aproximações gaussianas para fundos estocásticos \autocite{franciolini2025}. A contribuição desta dissertação é combinar a resposta dispersiva e a helicidade zero vinculada com comparações pareadas que distinguem compressão, aproximação probabilística e mudança de resposta. A geometria com posições de catálogo amplia os cenários estudados.

\section{Limitações e aplicação observacional}

O experimento é sintético: as posições celestes da população ampliada vêm de catálogo, enquanto distâncias, sinais e ruídos são simulados. O modelo periódico de Fourier e a fixação de parâmetros auxiliares em parte das campanhas delimitam a aplicação dos resultados. Em observações reais, o ajuste do modelo de temporização, a amostragem irregular e a inferência do ruído modificam a covariância e podem gerar pseudocovariância e acoplamento entre frequências.

A incerteza numérica foi incluída nos testes, conservando também as realizações cujas integrais não foram suficientemente resolvidas. Essa escolha evita selecionar apenas casos de integração fácil, mas reduz a capacidade de decidir alguns contrastes. As tolerâncias foram avaliadas por refinamentos e integrações independentes nas configurações estudadas. As comparações inconclusivas indicam onde concentrar novos cálculos.

O experimento de timing da Seção~\ref{sec:timing-sintetico} explicita esse mecanismo: a remoção de termos polinomiais e anuais gera pseudocovariância e mistura de canais mesmo com cadência regular. Na configuração examinada, ignorar esses dois efeitos produz uma divergência de 11,22 nat entre as distribuições dos dados. Esse valor caracteriza o exemplo de um pulsar; a campanha de inferência de massa mantém seu modelo periódico original.

A passagem para uma aplicação observacional envolve reproduzir uma análise pública de referência, incorporar a projeção do ajuste de timing à covariância da rede e inferir conjuntamente os parâmetros de ruído. Essa sequência permitirá avaliar os efeitos da aproximação dispersiva sobre restrições empíricas da massa e sobre buscas de polarizações adicionais.
